\subsection{Machine Learning and AI Research}

Machine-learning and AI claims combine statistical, computational, system, and social evidence, so the study design must identify which authority supports each conclusion.

\begin{process} \label{proc:ml-ai-study-protocol} \hlt{Machine-Learning and AI Study Protocol}\\
Design an empirical study before interpreting its scores.
\begin{enumerate}[label=(\roman*)]
\setlength{\itemsep}{0pt}
\item \hlt{Question}: state the population, task, context of use, intervention or comparison, and claimed outcome.
\item \hlt{Claim type}: classify each target as theoretical, empirical, causal, computational, system, or normative.
\item \hlt{Data}: record provenance, version, license, collection, consent or access status, labels, exclusions, and preprocessing.
\item \hlt{Evaluation boundary}: protect test data from training, model selection, prompt design, preprocessing, and human iteration.
\item \hlt{Comparators}: choose relevant baselines, ablations, and resource-matched alternatives before inspecting final results.
\item \hlt{Protocol}: fix splits, metrics, hyperparameter selection, seeds, stopping rules, hardware, software, and compute budget.
\item \hlt{Uncertainty}: repeat stochastic runs and report the source and calculation of intervals, variation, or significance measures.
\item \hlt{Stress tests}: examine subgroups, shift, robustness, failure modes, sensitivity, and plausible misuse within scope.
\item \hlt{Artifacts}: preserve code, data handles, checkpoints, configurations, logs, and command-to-result mappings.
\item \hlt{Report}: separate observed results from interpretation and state limitations, costs, risks, and unsupported generalizations.
\end{enumerate}
The NeurIPS checklist and its reproducibility program emphasize assumptions, complete proofs, experimental details, uncertainty, artifacts, ethics, and limitations \citep{neurips2026checklist,pineau2021reproducibility}.
\end{process}

\begin{process} \label{proc:ml-artifact-documentation} \hlt{Dataset and Model Documentation}\\
Release documentation with each durable ML artifact.
\begin{enumerate}[label=(\roman*)]
\setlength{\itemsep}{0pt}
\item \hlt{Dataset}: record motivation, composition, collection, preprocessing, labeling, intended uses, restrictions, and maintenance.
\item \hlt{Model}: record identity, version, intended and excluded uses, training data, evaluation protocol, subgroup results, and limitations.
\item \hlt{System}: record interfaces, human roles, deployment assumptions, monitoring, update policy, safeguards, and rollback conditions.
\end{enumerate}
Dataset datasheets and model cards establish complementary documentation patterns for provenance, intended use, evaluation, and limitations \citep{gebru2021datasheets,mitchell2019modelcards}.
\end{process}

\begin{principle} \label{princ:contextual-ai-risk} \hlt{Contextual AI Risk Assessment}\\
Accuracy alone is not a complete system claim. For the intended context, select and justify tests for validity and reliability, safety, security and resilience, accountability and transparency, explainability and interpretability, privacy, and harmful-bias management. These characteristics can conflict and require documented tradeoffs. This versioned vocabulary follows the voluntary NIST AI RMF 1.0 \citep{tabassi2023airmf}.
\end{principle}

\begin{warning} \label{warn:ml-evaluation-leakage} \hlt{Adaptive Evaluation and Leakage}\\
Any choice influenced by protected evaluation data makes the final score adaptive. Repeated leaderboard queries, test-set debugging, manual example selection, and unreported prompt tuning can all leak evaluation information. Reserve a final assessment boundary or disclose the adaptation and narrow the claim.
\end{warning}

\begin{warning} \label{warn:ml-benchmark-overclaiming} \hlt{Benchmark Overclaiming}\\
A higher score supports a claim only under the reported dataset, metric, resources, selection procedure, and uncertainty. It does not by itself establish practical usefulness, causal benefit, robustness, fairness, safety, or superiority in another operating environment.
\end{warning}
